\documentclass[11pt]{article}

% =====================================================================
% Crypto Was All You Needed
% The signed sovereign computation stack — Unbrowse, ZK identity, fair settlement.
% =====================================================================

\usepackage[margin=1in]{geometry}
\usepackage{hyperref}
\usepackage{amsmath}
\usepackage{amssymb}
\usepackage{enumitem}
\usepackage{xcolor}
\hypersetup{colorlinks=true, linkcolor=blue, urlcolor=blue, citecolor=blue}

\title{\textbf{Crypto Was All You Needed}\\
\large A Signed, Layer-Descending Stack for Agent Computation:\\
Identity, Authentication, and Zero-Knowledge Privacy across the Stack}

\author{Lewis Tham\\ Unbrowse AI \\ \texttt{lewis@unbrowse.ai}
  \and Cayden Chik\\ Unbrowse AI \\ \texttt{cayden@unbrowse.ai}}
\date{May 2026}

\begin{document}
\maketitle

\begin{abstract}
An agent that acts on the web stops being a polite HTTP client the instant it has
to click something, run a shell command, open a real browser, or put bytes on a
socket a bot-detector is inspecting. The load-bearing work does not confine itself
to the layer where JSON comes out, and neither can the security model. This paper
makes one claim and tries not to oversell it: the discipline that makes a
\emph{single} signed request trustworthy --- sign it, seal it, cache it, prove who
you are without showing your hand --- is the same discipline, unchanged, all the
way down the stack (screen $\to$ browser $\to$ CLI $\to$ OS $\to$ kernel $\to$
packet).\footnote{The title nods to ``Attention Is All You Need''; the argument
does not depend on the joke.} Every layer is signed by one Ed25519 wallet key;
credentials are bound to that identity by zero-knowledge proof and revealed only
under signature; every intermediate result is content-addressed and cached, with
the signed commitment appended to a hash-chained ledger (value off-chain, root
on-chain); and no accepted transition crosses a layer boundary without a signed
commitment. We cite an established primitive for each layer rather than inventing
one; the contribution is their
\emph{composition} into a single signed identity that owns the whole descent. Route
discovery --- turning a per-call browser tax into a shared, reusable route graph
--- is an orthogonal result established separately~\cite{wedge}; this paper is the
security, authentication, and privacy of the stack that \emph{executes} those
routes. Compensation sits above the substrate and stays deliberately narrow:
discovery and internal-API routing are free, and paid execution settles over x402.
\end{abstract}

\section{Introduction: one layer is not the stack}

The interface an agent uses to read and act on the web resolves an intent to a
ranked endpoint shortlist and executes the chosen call against the live site. That
much can be made cheap and reusable, and is~\cite{wedge}. This paper is about the
security model for what happens \emph{when the chosen action cannot stay at the HTTP
layer} --- which, for any agent that does real work, is often.

The common assumption is that the load-bearing part of agent work lives at the HTTP
layer, where things are clean. They are not clean. A login form is a screen-layer
fact. A session cookie is a browser-layer fact. A keychain entry is an OS-layer
fact. A TLS fingerprint is a packet-layer fact a defender is keying on whether or
not the agent was thinking about it~\cite{ja3,ja4}. An agent that is genuinely
accountable for its own actions has to be the same coherent, signed entity across
all of those layers --- not just the one that happens to return
\texttt{application/json}. A signature that secures the JSON call but not the
keychain read, the browser session, or the packet on the wire secures nothing the
adversary cannot route around.

So the thesis fits on one line: \emph{the layering is the point, and one signing
discipline holds at every layer.} The end-to-end argument in system
design~\cite{e2e} already established where each function belongs; we take it at its
word and run it down a signed agent stack. The result is a single Ed25519 identity
that owns every layer of the descent, credentials bound to it without disclosure,
and a content-addressed, sealed, ledgered substrate underneath --- each layer
pinned to an existing cryptographic primitive rather than a new one.


\section{Related work}

Our contribution is not a new agent or a new chain; it is the claim that one signed
discipline holds across every layer an agent touches, and that the route economy is
its natural settlement. We situate that against four lines of prior work.

\paragraph{Web agents and the browser-first assumption.} A large body of work has
agents drive real browsers: ReAct interleaves reasoning and action~\cite{react},
and benchmarks such as Mind2Web~\cite{mind2web}, WebArena~\cite{webarena}, and the
GPT-4V grounding study SeeAct~\cite{seeact} measure agents on live or simulated
sites. This work treats the rendered page as the interface. We take the opposite
stance: the browser is a fallback,
and the durable interface is the first-party API the page already calls. Our stack
keeps the browser as the bottom of a descent, not the default.

\paragraph{Tool and API invocation.} Toolformer teaches a model \emph{when} to call
an API~\cite{toolformer}; Gorilla teaches it to emit \emph{correct} calls across
thousands of APIs~\cite{gorilla}; the Model Context Protocol standardises how a model
connects to tools and data~\cite{mcp}. These solve single-agent tool use. They do not
address what happens when many agents discover the same interface repeatedly, who
maintains that knowledge as sites drift, or how the cost of discovery is shared --- the
maintenance-and-economics problem we leave to future work.

\paragraph{Identity, transparency, and capabilities.} Our trust layer reuses
established cryptographic systems rather than inventing them: EdDSA
signatures~\cite{rfc8032}, zero-knowledge credential binding~\cite{zklogin,camlys},
Certificate-Transparency-style append-only logs~\cite{rfc6962}, the object-capability
model~\cite{ocap}, and the Dolev--Yao adversary~\cite{dolevyao}. The novelty is not the
primitives but their \emph{composition} into one identity that signs every layer and
one ledger that records every claim, with ERC-8004~\cite{erc8004} as the cross-agent
identity surface.

\paragraph{Agent payments and incentives.} x402 revives HTTP 402 as a real
micropayment rail~\cite{x402}. Prior agent-payment work largely stops at ``the agent
can pay.'' We extend it to ``the payment routes to everyone who created the value'' ---
indexer, domain owner, platform --- as a fair three-way settlement over x402, while
discovery and internal-API routing stay free. To our
knowledge the combination --- signed multi-layer descent, a content-addressed-plus-
ledgered substrate, and a contribution-weighted three-party settlement over x402 --- is
not present in prior work as a single system.

\section{The stack is one node, instantiated at every layer}

We model the stack as a layered tree in the OSI tradition~\cite{osi}, but oriented
around \emph{who is acting} rather than \emph{which bytes move}:
\[
\text{screen-clicks} \to \text{browser} \to \text{CLI} \to \text{OS} \to
\text{kernel} \to \text{packet}.
\]
These are not six layers held together by prose. They share one structure: each
layer is described by the same four fields, so we define the structure once and
instantiate it per layer. A layer node has exactly four fields:
\begin{itemize}[leftmargin=1.4em,itemsep=1pt]
  \item \textbf{subject} --- the action this node performs at its altitude.
  \item \textbf{verb} --- one of \emph{observe}, \emph{act}, \emph{verify}: the kind of step.
  \item \textbf{witness} --- a check returning settled / not (two independent
        corroborations, or the step is not accepted).
  \item \textbf{parent} --- the hash of the node one layer up: the edge that links
        a layer to the one above it.
\end{itemize}
The descent is then a uniform recursion over one schema: each layer is a node, and
its child is the same node one altitude lower. Using a single schema is the whole
point --- it is what lets \emph{one} signature discipline and \emph{one} cache
serve the entire stack instead of six bespoke integrations.

The end-to-end argument~\cite{e2e} supplies the rule for which node acts: a function
lives at the highest layer that can do it completely and correctly, and lower layers
are fallback, not first resort. A cached signed plan beats a CLI call beats a browser
action beats moving a mouse across pixels like it is 2004. But when the higher layer
simply cannot act --- no API, hostile anti-bot, a page that is 100\% JavaScript and
0\% mercy --- the agent descends to the child node, and at the bottom it emits packets
whose TLS/HTTP fingerprint is indistinguishable from a real browser via
curl-impersonate~\cite{curlimpersonate}, matching exactly the JA3/JA4 signature the
other side is keying on~\cite{ja3,ja4}. The browser, CLI, and HTTP layers and the
fingerprint-faithful fetch run in the product, and the uniform \emph{signed descent through every
layer} (screen $\to$ browser $\to$ CLI $\to$ OS $\to$ kernel $\to$ packet) ships in
production: \texttt{src/values/signed-descent.ts} signs one wallet root once and
threads a hash-chained, per-layer signature down the whole stack, with tests that any
tampered or reordered layer fails to verify. Ownership is vertical and rides the
\textbf{parent} edge: \texttt{src/values/wallet-hierarchy.ts} derives each layer's
wallet from its parent's (HKDF) and the parent signs the child's key, and
\texttt{src/values/layer-wallet-descent.ts} has each layer sign with its own
parent-owned wallet --- so the root wallet owns the screen wallet owns the browser
wallet, all the way down to the packet wallet. What remains referenced rather than
shipped is emitting \emph{real} signed OS/kernel/packet operations across platforms;
the signed descent record, the wallet hierarchy, and the fingerprint-faithful HTTP
emission are in the product today.

The reason any of this composes is that the structure at every layer is identical
--- observe, act, verify, under one signature --- so the system's mechanisms are not
a grab-bag of separate subsystems but the recurring concerns of that one structure,
each met by an existing primitive (Table~\ref{tab:atoms}). We did not invent these
primitives; we mapped each layer's concern onto the appropriate one.

\begin{table}[h]
\centering
\small
\begin{tabular}{@{}lll@{}}
\hline
\textbf{concern} & \textbf{realised as} & \textbf{primitive (cited)} \\
\hline
descent      & spawn a child node one layer down      & curl-impersonate / JA3, JA4~\cite{curlimpersonate,ja3,ja4} \\
the tree     & one node whose children are nodes      & end-to-end argument / OSI~\cite{e2e,osi} \\
identity     & the root node's key, inherited down    & ERC-8004 trustless agent~\cite{erc8004} \\
witness      & \texttt{node.witness} --- the settle check & zkLogin / Pedersen commit~\cite{zklogin,pedersen} \\
cache        & a node memoising its child             & Merkle content-addressing~\cite{merkle} \\
\quad $\hookrightarrow$ KV parallel & same node, token altitude & attention KV cache / PagedAttention~\cite{attention,pagedattention} \\
seal         & the gate on \texttt{node.verb}         & AEAD / object-capabilities~\cite{rfc5116,ocap} \\
ledger       & the \texttt{node.parent} edge, append-only & Certificate Transparency / hash chain~\cite{rfc6962,bitcoin} \\
loop         & a node re-firing until its witness settles & OODA / proof-of-history clock~\cite{ooda,solana} \\
\emph{(edge)} economics & payment on the parent$\to$child call & x402 split settlement~\cite{x402} \\
\hline
\end{tabular}
\caption{The system's mechanisms are the recurring concerns of one layer structure,
each aligned to an existing, cited primitive. One element is \emph{not} a node and we
name it rather than hide it: economics (x402) is a weight on the \textbf{parent}
edge, not a node --- so the precise structure is one node type plus one typed edge
type, a weighted tree (this is the whole content of \S\ref{sec:economics}). A
runnable reference implementation of the cache + ledger core ships alongside this
paper (\texttt{reference/}, with tests that execute each claim).}
\label{tab:atoms}
\end{table}

\subsection{Traversing the tree: descent is a cheapest-first search}

A tree of nodes is inert until something walks it. The stack's operations ---
resolve an intent, execute a route, fall back to a browser --- are not three
subsystems but \emph{one graph traversal of the node tree, weighted by compute
cost}. The end-to-end rule above is exactly Dijkstra's algorithm with edge weight =
compute cost: visit nodes in increasing cost and stop at the first whose witness
settles.
\begin{verbatim}
walk(intent):
  for node in tree.nodes_by_increasing_cost(intent):   # Dijkstra over compute cost
    if node.cached and node.witness.settled:            # RESOLVE -- cache hit, 0 descent
      return node.replay()                              # EXECUTE -- no deeper layer touched
    if node.can_act_completely(intent):                 # highest capable layer
      r = node.act()
      if r.witness.settled:
        cache.put(node, r); ledger.append(node)         # CAPTURE the win for next time
        return r
  return walk_child(intent)                             # FALLBACK -- the browser is the deepest child
\end{verbatim}
BFS, DFS, and Dijkstra are not three different walks here but \emph{one walk with
three priority functions} over the same tree: BFS = uniform cost, DFS = descend
immediately on miss, Dijkstra = cost-weighted (the production default,
cheapest-capable-first). The resolve$\to$execute$\to$fallback ladder \emph{is} that
Dijkstra instance --- three short-circuit tiers, cheapest first: cache replay (0
descent) $\to$ one capable-layer act $\to$ descend to the child. Each settled node
pushes future identical intents to a cheaper tier, so the deep layers (the browser)
are the cache-miss oracle, visited as rarely as possible.

\subsection{The traversal record is the ledger}

Every node the walk visits, and the verdict its witness returned, is appended to one
append-only log keyed by the node's content hash. That single object wears three
names and is the same primitive each time: as the \emph{route ledger} it records
\texttt{\{signer, value-hash, timestamp, prev-hash\}} per settled node (\S8); as the
\emph{resolution ledger} (\texttt{src/values/resolution-ledger.ts}) it is the
content-addressed memo of the traversal, so a hit short-circuits the whole descent
(the cache of \S6); and as a \emph{session record} it is the human-readable tail of
the same log, where a node recorded \emph{refuted} is a pruned branch the walk never
re-enters --- the negative result is a cache entry, not a deleted fact. Memoisation
(do not recompute a settled node), pruning (do not re-walk a refuted node), and
provenance (the hash-chain proves the order) are one append-only ledger seen three
ways. \S8 is that ledger as the \textbf{parent} field.

\subsection{Composition: a walked traversal is itself a node}

Descent makes a node's \emph{child}; the same schema run upward makes a node's
\emph{parent}. When one intent cannot be settled by a single route --- the route it
resolves to needs a binding only another route yields --- the walk visits several
nodes in dependency order, threading each one's output into the next. That ordered
sub-walk is not a transient. It is the traversal record of the previous subsection,
and by the same content-addressing it is itself a node: a \emph{composite} is the
parent whose children are the routes it composed, keyed by the content that makes it
the same DAG --- its domain, its target, its ordered constituent nodes, and the
binding edges between them --- and replayed as one unit. The first agent to walk the
chain records the composite; a later agent, including one that never walked it, looks
it up by that address and replays the recorded order rather than re-deriving it ---
exactly the cache-hit short-circuit of the walk (\texttt{node.cached}
$\Rightarrow$ \texttt{node.replay()}) one altitude up. So the recursion closes in
both directions: descent spawns the child when the parent cannot act, and composition
records the parent when several children together settle an intent no child settles
alone --- and the traversal record is the ledger at every altitude. This ships:
\texttt{src/orchestrator/} content-addresses a walked multi-step traversal as a
composite and, on a later resolve, replays the recorded order before re-walking, with
a constituent-staleness guard --- a removed or disabled child invalidates the
composite back to a full re-walk, so a stale parent never replays a broken child.

\section{Identity: the root node's key, inherited by every child}
Identity is not a layer of its own --- it is the one \textbf{subject} field every
node shares, and the property that the whole tree is keyed by one keypair. The root
of trust is a single Ed25519 keypair~\cite{rfc8032} --- and conveniently, the agent
already has one, because a Solana account literally \emph{is} a base58 Ed25519 public
key. The wallet is the identity. Every action --- a click, a syscall, a packet --- is
an admissible \texttt{node.verb} only if it chains to a signature under that one root.
No signature, no action; the stack has no anonymous side door. The Unbrowse signer
already produces wallet signatures for resolve/execute admission. The key descends the
\textbf{parent} edge, never re-keying: each child node's key is HKDF-derived from its
parent's and the parent signs it (\S3), so one identity threads the whole tree.

For trust \emph{between} agents we adopt the ERC-8004 ``Trustless Agents''
model~\cite{erc8004}, which defines three on-chain registries: \textbf{Identity}
(portable agent IDs), \textbf{Reputation} (signed feedback), and
\textbf{Validation} (independent re-execution / ZK / TEE checks). An Unbrowse route
maintainer is precisely one of these agents: a portable identity that can carry
reputation and be checked by someone who does not trust it yet. The three ERC-8004
record types map cleanly onto the wallet identity already in play --- a portable
Identity that is the wallet pubkey, signed Reputation feedback, and an independent
Validation re-execution record, each signed by a real wallet.
Binding to the \emph{deployed} on-chain registries remains integration work, and we
will not pretend otherwise.

\paragraph{Key mobility and the public boundary.} One root keypair owning every
layer raises the obvious question: what may be exposed, and where? The answer is an
asymmetry. The \emph{private} key is mobile \emph{inward} --- it descends or ascends
the stack to surface an identity or value at whatever altitude needs it, a click at
the screen layer or a signature on a packet, always the same root. Across the
\emph{public} boundary the rule inverts:
\textbf{only value copies cross the public boundary} --- content-addressed copies
of the value, verifiable against the
\emph{public} key and nothing more. The private key never leaves; what the world
receives is a copy it can check, but cannot forge and cannot reverse into the
secret. Surfacing a value at a layer is a private-side act under the one key;
publishing it emits a fresh, content-addressed copy carrying the public key and a
signature, never any private material. This is the runnable property in
\texttt{paper/reference/layers/key\_mobility.py}: one key surfaces a value at every
layer, a foreign wallet cannot forge a public copy, a tampered copy fails its
content hash, and the private key is provably absent from everything that crosses.
The production port ships in \texttt{src/values/wallet-seal.ts}: \texttt{publishValue}
emits a \textbf{wallet-signed value copy} verifiable under the public key, the
outward complement of the sealed-unless-revealed cache, with the private key loaded,
used once, and zeroed before the copy ever leaves.

\section{Witness without disclosure: ZK credential binding}

This section is \texttt{node.witness} made non-disclosing: the witness must settle
``yes, bound'' without revealing what it witnessed. Signing is the easy half. The
hard half --- the actual research --- is proving a credential belongs to the identity
\emph{without revealing the credential}. Your
password, your session cookie, your 1Password entry, the fact that you clicked
``approve'' for this domain: all of it should be provably bound to the wallet and
leak \emph{nothing} identifiable at any layer unless you choose to open the proof
under signature.

This is the classical anonymous-credential problem~\cite{camlys}, and it has grown
up. zkLogin~\cite{zklogin} binds an existing OAuth/OpenID identity to a chain
address through a Groth16 proof, hiding the link even from the identity provider;
Semaphore~\cite{semaphore} proves anonymous group membership with zk-SNARKs;
Pedersen commitments~\cite{pedersen} give the ``commit now, open later, can't lie
about it'' primitive for a single secret. Together they make the binding a
two-witness corroboration that betrays nothing: the chain sees \emph{that} a
credential is bound, never \emph{what} it is. This is the central
contribution, and the primitive is implemented in the product: \texttt{src/values/zk-binding.ts}
implements the same non-interactive Schnorr proof (Fiat--Shamir over a 2048-bit MODP
group) that proves a credential is \emph{bound to the wallet without revealing} it ---
the wallet signs $y=g^{x}$ where $x$ is derived from the credential, and a holder
proves knowledge of $x$ while the verifier learns only ``yes, bound.'' It is wired at
the capture boundary by \texttt{src/capture/zk-bound-hole.ts}: each redacted secret
\emph{hole} carries the binding, so the backend confirms a credential is bound to the
wallet without ever seeing the secret, and the holder proves knowledge only at fill
time. Tests confirm the credential never appears in the binding or the proof, that a
wrong credential or a foreign wallet cannot forge one, and that an unbound hole fails
closed. The remaining integration is adopting this binding across every entry of the
live credential vault end-to-end; the primitive and its hole-level binding are shipped
and tested, not promised.

\section{Cache: a node memoising its child (content-addressed, sealed unless revealed)}

The cache is simply a node that memoises its child: address the result by the hash of
its own bytes, and a later walk that reaches the same node replays the stored value
instead of descending again (the traversal short-circuit of \S3.1). Every layer's
settled result --- a resolved endpoint, a rendered page, a signed plan --- is memoised
under a content-addressed key, in the Merkle
tradition~\cite{merkle} and exactly as realised by content-addressed stores like
IPFS~\cite{ipfs} and Git's object model~\cite{gitobjects}. Same structure
re-walked, same key, same answer, deterministically; the cache is re-derived to
\emph{verify}, never trusted because it looked right last time. Unbrowse's
centralised cache servers can hold these entries and serve them across agents,
while a sensitive entry is kept sealed until the holder opens it --- a
content-addressed cache that is public by default and private by proof. Endpoint and
route caching run in the product, and the sealed-unless-revealed commitment layer now ships in
production: \texttt{src/values/wallet-seal.ts} addresses each value by the
sha256 of its \emph{plaintext} (so the same content resolves to the same key on any
host) yet stores AES-256-GCM ciphertext under a key bound to the wallet, and tests
confirm the at-rest bytes are unreadable, that only the binding wallet can reveal, and
that a tampered ciphertext refuses to open.

\subsection{The route graph as a commitment-only structure}

The content-addressed cache above keys each value by the sha256 of its plaintext, so the
same content resolves to the same key on any host. That is exactly right for a
\emph{private} cache --- but the route graph is \emph{shared} across agents, and a bare
hash on a shared structure leaks: a low-entropy value (a boolean, a locale, a short id) is
recoverable by hashing a dictionary against the visible commitment, and equal values expose
linkable equality across hosts. So a captured credential or storage value --- a session
token, an anti-bot \texttt{localStorage} entry --- is promoted to a first-class node in the
requires/yields graph under a keyed commitment: a MAC under a wallet-derived key, not a bare
hash. An observer of the shared graph who lacks the wallet cannot brute-force a low-entropy
value, and the same value under two different wallets yields two different commitments, so
nothing links across installs. The value itself is sealed off-graph under the wallet and
revealed only locally, per step, at the moment of replay --- the graph carries the
dependency \emph{shape} (which operation yields a token, which operation requires it) while
the secret is held by no one but the holder. The keyed commitment and its graph binding are
implemented in \texttt{src/values/storage-hole-bindings.ts} over the off-graph store
\texttt{src/values/sealed-blob-store.ts}, and revealed at the replay boundary in
\texttt{src/execution/index.ts}; tests confirm the secret appears on neither the graph node
nor the bytes at rest, that a wrong key cannot reveal, and that a legacy plaintext token is
never mistaken for a commitment.

The cache and the ledger meet in one production primitive worth naming, because it is
how a slow truth-resolution is paid for once and never again.
\texttt{src/values/resolution-ledger.ts} treats a signature as the KV key to a
\emph{pointer} that resolves on truth resolution. An intent is content-addressed to a
pointer; the first resolution runs the expensive work, stores its result content-addressed
(\texttt{sha256:}$\langle$hex$\rangle$, exactly the \texttt{putBlob}/\texttt{resolvePointer}
shape), and appends one row to an append-only, hash-chained \emph{ledger of resolutions}.
Every later call resolves the pointer instead of recomputing --- and because the value is
addressed by the hash of its own bytes, an evicted entry rebuilds the \emph{identical}
layer, exactly like a Docker layer cache: a miss rebuilds it, a hit replays it.
A settled value is a \emph{promise} in the two senses computer science already gives the
word: a \emph{future}~\cite{futures} --- a placeholder computed once and thereafter only
read, never recomputed --- and a promise in Burgess's \emph{Promise Theory}~\cite{promisetheory},
where an autonomous agent can promise only its \emph{own} behaviour and the order of the
whole is nothing but the body of promises voluntarily kept. The resolution ledger is exactly
that body: each layer, as an autonomous agent, promises the single value it derived, addressed
by the hash of that value's own bytes --- so it is a promise the agent can always keep (derived
once, then \emph{true forever}), and the append-only, hash-chained ledger is the public,
independently corroborable record of those promises. Because it
is addressed by the hash of its own bytes, the same value resolves to the same pointer
in every process, so a memo keyed on that pointer hits forever and never recomputes ---
asking ``does it need to re-resolve each time?'' answers itself: no, not if it is true
forever. The only thing that busts the cache is the \emph{value} changing, and the
sharpest case is when \emph{time is part of the value}: then the pointer changes every
moment and the derivation re-resolves each time --- unless the promise included time
itself, a timeless claim is computed once and held. Invalidation is this same
discipline run the other way: when an upstream value changes, its bytes change, its
content hash changes, the pointer keyed on the old hash no longer resolves, and the
layer re-resolves on next access (\texttt{src/values/resolution-ledger.ts}). Nothing is
invalidated by a clock or a manual flush; correctness falls out of the addressing,
because a changed input is \emph{by construction} a different key --- the same reason
changing a base layer in a Dockerfile invalidates the layers built on it. Crucially,
dependent recompute needs \emph{no} dependency-graph walk: a value keyed on another
value's pointer is automatically a different key the instant that pointer changes, so a
changed input invalidates everything addressed through it without a cascade to engineer.
The fallback that feeds it walks the descent ladder highest-capable-first
(\texttt{src/values/kv-fallback-pipe.ts}), content-addressing each layer's result, so a
hit short-circuits the whole descent and a changed input re-resolves only at the layer
it touched. The recompute boundary is proven runnably
(\texttt{reference/ledger/recompute.py}): a timeless value derives once over a hundred
reads, a time-keyed value recomputes each moment, and a dependent value re-derives
automatically when its upstream pointer flips. Tests confirm the warm path never
recomputes, the ledger is tamper-evident, and eviction reproduces the same pointer.


The parallel worth naming is to the \emph{KV cache} inside the transformer doing
the reasoning. Attention~\cite{attention} computes a key and value for every token;
the KV cache memoises those states so decoding never recomputes them, and modern
serving engines page that cache like OS virtual memory and \emph{share identical
prefixes across requests} by content (hash) key~\cite{pagedattention}. That is the
same discipline, one altitude up: same prefix, same cached blocks, reused not
recomputed --- except the structural key is now a stack subtree rather than a
token prefix, and the entry is ZK-sealed rather than served in the clear. The model
caches computation it has already done; the stack caches \emph{actions} it has
already signed. Same shape, different altitude --- a KV cache all the way down,
where each layer's entry stays sealed until a signature opens it.
\section{Seal: the gate on \texttt{node.verb} (nothing ships unsigned)}

The seal is not a layer either --- it is the gate every node's \texttt{verb} must pass
before it emits, and authority descends the \textbf{parent} edge attenuated, never
widened. Before any layer emits, it self-verifies through the root. At the wire that is
authenticated encryption with associated data (AEAD)~\cite{rfc5116}: the receiver
checks the tag before it trusts a single byte, and tampering fails closed rather
than failing quietly into your logs. The same discipline generalises upward as
object-capability attenuation~\cite{ocap} --- authority travels only by reference
and can be narrowed, never picked up from the ambient air. The rule is boring on
purpose: no unsigned action crosses the gate, at any layer, ever. This holds as a
\emph{uniform} stack property: \texttt{reference/layers/gate.py} is a runnable gate
bound to one wallet root that admits an action only if its signature verifies over
the action's canonical bytes, and a test confirms that an unsigned, foreign-signed,
or tampered action is rejected and counted --- zero unsigned actions ever cross.
AEAD itself is, obviously, ubiquitous and shipped at the transport layer; we are not
claiming to have invented encryption.

\section{Ledger: the \texttt{node.parent} edge made append-only}

Stated as a node field, the ledger is what happens when \texttt{node.parent} is
required to be append-only and ordered (\S3.2): the \texttt{prev-hash} \emph{is} the
parent edge, realised as a chain rather than a free pointer. Signing is only half the
story; the obvious next question --- and the one most ``signed'' systems wave away ---
is \emph{where does the signature go}. A cache and
a ledger are not the same object, and conflating them is the bug. A cache is
content-addressed: you fetch a value by the hash of its content and order does not
matter, so the same hash resolves on any host~\cite{merkle}. A ledger is the
orthogonal half: append-only, \emph{ordered}, and hash-chained, so the history and
the order are themselves tamper-evident. A signed entry --- a signed KV result, a
route attestation, a signed quality claim --- needs both: the cache holds the
value, the ledger holds the signed commitment to it.

The load-bearing design choice is \emph{value off-chain, root on-chain}. You never
put the payload on a chain --- that would be paying gas to store megabytes and
reading them slowly. You store the small, signed commitment: an entry of
\texttt{\{signer, value-hash, timestamp, prev-hash\}}, where the value itself lives
in the content-addressed cache and only its hash appears in the log. Each entry
chains to the prior, so reordering or editing any past row breaks every root after
it~\cite{bitcoin}. Independent auditors --- not a trusted operator --- verify the
log is append-only and consistent, exactly the model Certificate Transparency
formalises with its Merkle log, Signed Tree Head, and inclusion
proofs~\cite{rfc6962}. That is what makes the ledger trustless rather than
``trust our server.''

This also answers how the system decentralises without a rewrite. Per-entry
on-chain writes do not scale, so you batch: thousands of entries collapse to one
Merkle root, and only the \emph{root} is checkpointed on-chain (the
Certificate-Transparency / rollup pattern), ordered by a decentralised clock such
as Solana's proof of history~\cite{solana}. The maturity ladder is one shape at
every rung: a local append-only log today, a signed-root server next, on-chain
Merkle-root checkpoints as the decentralised end-state --- you swap the host, the
signatures and the hash-chain never change. The local append-only
ledger runs in the product (the route graph already keeps signed JSONL records). The
Merkle-root checkpoint and independent-auditor layer ship as runnable reference code:
\texttt{reference/ledger/checkpoint.py} batches many signed entries into one root
and emits a per-entry inclusion proof any auditor verifies against that root, with a
test that an entry outside the batch cannot prove inclusion and that changing any
entry changes the root. Publishing the root on-chain is the deployment step; the
commitment it would publish is exactly this root.

\section{From security to economics: where this paper stops}
\label{sec:economics}

A signed, sealed, ledgered route graph is a security substrate; what it is
\emph{worth} --- who maintains the routes, who pays, who earns, and how a freshness
claim is made costly to fake --- is a separate concern with its own
paper~\cite{maintenance}. We draw the line here deliberately. This paper is the security, authentication, and privacy
of the stack: one key signs every layer, credentials are bound by zero-knowledge
proof and revealed only under signature, every result is content-addressed and
sealed, and nothing ships unsigned. Fair compensation that sits on top of that
substrate is deliberately simple: discovery and internal-API routing are free, and
paid execution settles fairly over x402 across the parties who created the value.

Keeping security and compensation as separate concerns is itself the discipline. A security
model must stand on its own, without leaning on a token to be true: the signatures,
the ZK binding, and the sealed cache are correct or not on cryptographic grounds
alone, independent of any economics layered above them. Conversely, the economics
is only worth stating once the substrate it secures is trustworthy. So we settle
the substrate here; the compensation rule above is the whole of what sits on top of
it --- free discovery and routing, paid execution settled fairly over x402.
\section{The control loop}

The whole stack runs as an OODA loop~\cite{ooda}: observe the layer's state, orient
it against the model (this is where the learning actually lives, despite ``decide''
getting all the press), decide the next action at the right layer, act under
signature --- then repeat with feedback, dropping a layer or replanning on failure,
and stopping when the goal is reached. It is the same observe--orient--decide--act
cycle run uniformly at every layer of the stack. If that sounds unglamorous, good:
the glamorous control loops are the ones that do not converge.



\section{Evaluation: the security substrate}

This paper's claims are about a security discipline, so we evaluate the security
substrate it introduces --- not the end-to-end latency win of route reuse, which is
a discovery result established separately. We
are precise about what the measurements here do and do not establish.

\paragraph{The substrate is correct by executable test.} The local cache+ledger
reference implementation that ships with this paper is unit-verified:
content-addressing, value-off-chain/root-on-chain separation, hash-chain
tamper-evidence, ed25519 signatures, and deterministic Merkle-root inclusion all
pass as executable tests, so the substrate's correctness claims are runnable, not
asserted.

\paragraph{Sealed-cache reuse, measured.} We additionally ship a runnable
micro-benchmark (\texttt{reference/bench/\allowbreak bench\_reuse.py}) that isolates the one
property the sealed cache rests on --- content-addressed reuse versus re-derivation
--- on the actual cache implementation. Over many trials, re-deriving a representative
$\sim$64\,KB extracted payload costs milliseconds while a content-addressed cache hit
(read + hash re-verify) costs tens of microseconds: a large mean speedup, wall-clock,
that lands in the \textbf{$\approx$50--90$\times$} range hardware-dependent, with a gate
that fails (and is the part actually pinned) if reuse is not materially faster
($>2\times$). This is a CPU-only
recompute-versus-reuse demonstration --- it isolates the cache primitive, not the
network or the browser --- and it establishes exactly one thing: a sealed,
content-addressed commitment is asymptotically cheaper to re-read than to
re-derive, which is what makes the cache safe to lean on rather than merely fast.

\paragraph{The descent reaches the network interface.} The security
discipline is only as deep as its lowest layer, so we are concrete about the bottom
of the descent. The fingerprint-faithful fetch is not only a diagram: it is
implemented as the
orchestrator's curl-impersonate fetch (\texttt{src/capture/curl-impersonate-fallback.ts})
backed by a vendored uTLS CONNECT-proxy daemon (\texttt{src/cdp/proxy/utls-daemon.ts})
across four platforms (darwin/linux $\times$ amd64/arm64), so the agent's TLS
ClientHello and HTTP/2 settings reproduce a real browser's JA3/JA4 signature at the
\emph{network interface}, not merely a spoofed User-Agent header.
This is the packet layer of
screen$\to$browser$\to$CLI$\to$OS$\to$kernel$\to$packet made real: the same signed
identity that authorises a route also emits the bytes that carry it, and the bytes
are browser-indistinguishable on the wire.

\paragraph{Coverage of the full descent, measured.} Because the descent is what lets
the agent reach data behind TLS-fingerprinting anti-bot, we measure whether it
actually returns results across a broad, diverse intent space rather than on
cherry-picked domains. A live coverage gate runs the real \texttt{unbrowse search}
binary --- which exercises the whole descent including the packet-layer fetch ---
over a diverse intent set and counts the fraction that return real results; on the
current graph it returns results on every probed intent (coverage $=1.0$ against a
$\geq 0.75$ release threshold). \texttt{scripts/coverage-gate.sh} is
re-runnable and fails honestly when the descent cannot reach the data. The dollar
saving the wedge measures~\cite{wedge} --- a cold browser rediscovery at
\$0.10--0.53 collapsed to a \$0.005--0.02 one-time install --- rides on this same
descent: cheaper because the packets are emitted once and the route reused, not
re-derived through a browser on every call.

\paragraph{Per-task retrieval behind anti-bot, measured.} We now measure directly
what an earlier draft deferred: per-task retrieval on an adversarial,
heavily-JavaScript-gated site. On a nine-post corpus drawn from three communities of
a major social platform whose HTML surface is JavaScript-challenge-gated and whose
read endpoint returns HTTP~403 to a naive client --- ground-truthed against the
platform's own listing data --- a naive HTTP client is blocked on \textbf{100\%} of
requests, while the descent retrieves the real content on \textbf{9/9} posts
(400--820\,KB each) and recovers the ground-truth author handle and a distinctive
title token on \textbf{100\%} of them. This \emph{anti-bot retrieval} head-to-head
(a re-runnable anti-bot retrieval suite) is the per-task adversarial-site measurement the
descent is built for: naive 0/9, descent 9/9, on a site that 403s ordinary scrapers.

\paragraph{Execute, don't guess --- the same principle, measured at model scale.} The
discipline this paper applies to the web --- call the real interface and execute it,
rather than have an agent re-derive it --- holds for models too: route to a real tool
and execute, instead of guessing from weights. A reproducible, gated benchmark suite
(\texttt{bench/BENCHMARKS.md}) shows a small on-device model (Qwen2.5-1.5B) routed to a
library of executable tools turning tasks it fails from weights alone into tasks it
solves reliably. Two are genuine tools-versus-no-tools gains on the same 1.5B model:
knowledge absent from the weights \textbf{0\% $\to$ 95\%} by retrieve-then-execute, and
applying a retrieved skill rather than reasoning it from scratch \textbf{63\% $\to$ 93\%}.
Two are distillation gains on the served model: code-correctness \textbf{68\% $\to$ 100\%}
(raw base to distilled, in-distribution) by routing to a real executor, and hard-reasoning
families \textbf{50\% $\to$ 92\%} by distilled routing --- measured against a
trained-specialist baseline, a scope we flag rather than overclaim. The architecture is the
capability, not the raw weights --- the same claim the route graph makes for the web.

\paragraph{Self-improving by reuse.} The system improves by running against itself.
Resolving a fixed probe set repeatedly, latency falls \textbf{21.1\,s cold $\to$ 4.1\,s
warm (\textbf{$-$80.7\%})} as the route cache fills, then plateaus (tail spread 4.9\%
over the last five of twenty passes). The plateau is a physical limit, not a tuning
choice: once every route is cached, further passes cannot reduce latency. Reuse, not
retraining, is where the compounding comes from.

\paragraph{Credentials are wallet-bound, witnessed end-to-end.} The auth descent is
exercised against live endpoints: a credential is \emph{sealed to the holder's wallet},
revealed only for the correct key (a wrong key fails closed), and attached to a real
authorised request --- an authenticated search and an authenticated action each
round-trip successfully, with the agent ever holding only a content-addressed
commitment, never the raw secret (\texttt{bench/agent-experience/}, re-runnable). Sealing,
fail-closed reveal, and the authorised round-trip are all witnessed, not asserted.

\paragraph{The pointer-only invariant, scanned.} The claim that the stateless
binary writes no secret to disk --- only pointers and signatures cross the wire ---
is not left as an assertion. A re-runnable capability gate
(\texttt{bench/capability/}) scans every persisted session file and counts any
plaintext secret. On the current build it scans \textbf{446} session files and
finds \textbf{0} leaks: the redaction invariant (sealed value, pointer at rest)
holds across the whole on-disk surface, not just the paths under test. This is the
sealed cache of \S6 measured as a property of the product, not only of the
reference implementation.

% BEGIN browsecomp-self-improvement (generated by inject-paper-sentence.py)
\paragraph{BrowseComp and warm-cache self-improvement, measured.} On OpenAI's
\textbf{BrowseComp} multi-hop browsing benchmark, driving the same agent and grader
through the route-graph search path across repeated tries, the route/content cache
warms run-over-run: per-query wall-clock falls from \textbf{80.8\,s} on the cold
graph to \textbf{65.4\,s} once warmed (\textbf{$-$19\%}), the
capture$\to$index$\to$reuse self-improvement the sealed cache predicts. Accuracy is
dominated by the agent harness above retrieval, not by the substrate, and sits
honestly below specialised search stacks: our reproducible single-shot figure is
\textbf{0.100} against Exa's published \textbf{0.336}, a gap we report rather than
hide. We record the full per-try ledger in
\texttt{bench/browsecomp/SELF-IMPROVEMENT.md} rather than headline a number this
minimal agent does not earn; where the substrate \emph{does} win head-to-head is
anti-bot retrieval, measured next.
% END browsecomp-self-improvement
\paragraph{What this establishes, and what it does not.} The substrate tests
establish that the security primitives are correct as implemented; the
micro-benchmark establishes that sealed-cache reuse is cheap; the vendored uTLS layer
and the coverage gate establish that the descent reaches the network interface and
returns real results broadly; the anti-bot and wallet-auth measurements establish that
the descent retrieves real content past JavaScript-challenge anti-bot and that
credentials bind to the wallet end-to-end. What remains out of this paper's scope is
\emph{end-to-end multi-hop task accuracy} on the open web --- which is dominated by the
agent harness above the retrieval layer, not the substrate (we report our reproducible
BrowseComp figure and its warm-cache behaviour in the repository, honestly below
specialised search stacks) --- and emitting \emph{real signed} OS/kernel/packet
syscalls across platforms: the fingerprint-faithful HTTP emission is shipped; raw
signed syscall descent is referenced, not claimed.
\section{Threat model}

We state the adversary explicitly, because a trust paper that never names what it
defends against is decoration. We assume a Dolev--Yao network
attacker~\cite{dolevyao} --- able to read, drop, replay, and forge messages on any
layer --- plus three application-level adversaries the route economy specifically
invites. For each, we name the atom that resists it and the residual risk we do not
claim to close.

\paragraph{(A1) The impersonator.} An attacker replays or forges an action to act
as another identity. Resisted at the \emph{root}: every action carries an Ed25519
signature over its canonical bytes~\cite{rfc8032}, and the ledger's
\texttt{prev}-hash makes a replayed entry detectable out of order~\cite{rfc6962}.
Residual: key theft is out of scope --- a stolen wallet key is the user, by
construction. This is the same boundary every signing system accepts.

\paragraph{(A2) The credential thief / over-reach.} An operator retargets imported
sessions at accounts they do not own --- the abuse vector that closed the
source~\cite{ossnotice}. The \emph{witness} atom narrows it: credentials are bound
to the identity by zero-knowledge proof and revealed only under
signature~\cite{zklogin,camlys}, so a leaked ledger or cache exposes \emph{that} a
credential is bound, never the credential. Capability attenuation~\cite{ocap} bounds
what a delegated action may do. Residual: an operator authenticated as themselves
can still drive their own credentials; we constrain reach, not self-harm.

\paragraph{(A3) The poisoner.} An adversary publishes a false or stale route to
mislead later agents --- the integrity attack on a shared graph. Resisted at
\emph{witness} and \emph{ledger}: every route attestation is signed and appended to
the hash-chained log~\cite{rfc6962}, and finalisation can require a $t$-of-$n$
quorum~\cite{frost}, so one dishonest maintainer cannot unilaterally settle a claim
and any tampered or reordered attestation is detectable. Residual: routes are
re-derived to verify rather than trusted on sight; the model moves trust-sensitive
traffic to corroborated routes, it does not claim every route is adversary-proof.

\paragraph{(A4) The free-rider.} An actor consumes paid execution without paying.
The \emph{verb} atom binds payment to the act: an unpaid \texttt{execute} gets HTTP
402, not service~\cite{x402}, while discovery and internal-API routing stay free.
Residual: the free tier is, by design, free to consume; the payment gate sits at
execution, which is where the cost actually lands.

\paragraph{(A5) The public-boundary leak.} An adversary --- or an honest
contributor by accident --- leaks the closed capture/integrity engine into a public
artifact. Resisted mechanically, not by policy: \texttt{leak-guard.sh} and
\texttt{paper-gate.sh} run in release CI and fail the build if a sensitive term or
an unanchored claim reaches a public path. This paper is itself subject to that
gate. Residual: NDA source review remains the only full-disclosure path, by the
abuse reasoning of~\cite{ossnotice}.

\paragraph{What we do not defend.} We make no claim against a malicious endpoint
that lies in its own responses (a route can faithfully replay a hostile API), a
global passive adversary correlating timing across the whole network, or coercion of
a key holder. These are out of scope and named so the reader is not misled by
silence.

\section{What is built, what is referenced (no fabricated green)}

In the spirit of not selling a roadmap as a changelog, we separate what runs in the
product, what is available as runnable code, and where the work stops:
\begin{itemize}[leftmargin=1.4em]
\item In the product: Intent $\to$ route resolve $\to$ execute; live browser capture; HTTP
fetch with browser-faithful TLS fingerprinting; wallet-signed admission;
route/endpoint caching.
\item The cross-layer security primitives are implemented, each with
tests that execute the claim: the signed descent through every layer with vertical
wallet ownership (\texttt{src/values/signed-descent.ts}, \texttt{src/values/wallet-hierarchy.ts},
\texttt{src/values/layer-wallet-descent.ts}), ZK credential binding --- the central
contribution --- and its capture-boundary wiring (\texttt{src/values/zk-binding.ts},
\texttt{src/capture/zk-bound-hole.ts}), the sealed-unless-revealed cache
(\texttt{src/values/wallet-seal.ts}), the signature-keyed, recomputable resolution
cache and its ledger of resolutions (\texttt{src/values/resolution-ledger.ts},
\texttt{src/values/kv-fallback-pipe.ts}), and the hash-chained signed ledger
(\texttt{src/values/sealed-ledger.ts}). The architecture is
backend-as-harness --- the client surfaces only the holes to fill plus wallet-sealed
auth, and the backend returns nothing but structure; the composed flow and its
sealing primitives are implemented server-side
(\texttt{src/capture/backend-reveng-endpoint.ts}), so the reverse-engineering engine
no longer rides in the client. The stateless binary writes no local state --- only
pointers and signatures cross the wire, a property measured by the 446-file leak
scan reported in the evaluation.
\item The runnable reference suite in \texttt{reference/} remains the
cited foundation for each primitive above and for the parts not yet in the product:
Merkle-root checkpoints with inclusion proofs (\texttt{ledger/checkpoint.py}) and the
inverse client/OS harness as a capability-gated, content-addressed pipe
(\texttt{pipes/pipe\_contract.py}). The companion Aiko engine mirrors this boundary:
the server owns recursive contract compilation, the CLI is the bridge, and the
client sees only holes, approvals, and typed pointers rather than graph-control
verbs. The
on-chain deployment of the checkpoint root, the ERC-8004 registry binding, and a
real signed OS/kernel/packet descent across platforms remain integration work,
honestly labelled.
\item Zero-edit drop-in replacements for the libraries an agent would
otherwise reach for exist as packages backed by the route graph: an \texttt{exa-py}
drop-in (\texttt{packages/py-exa}) and a \texttt{browser-use} drop-in
(\texttt{packages/py-browser-use}), each providing the upstream's public surface so a
single import swap routes the call through Unbrowse instead.
\end{itemize}
Treat the reference suite as running, tested code with cited foundations --- not a
feature list with a release date, and not a research agenda either. The gap between
running code and a press release is the entire difference between a whitepaper and a
pitch, and we would like to stay on the correct side of it.

\section{Conclusion}

An agent that is accountable for its own actions has to be one coherent entity
across every layer it touches: under one key, revealing nothing it did not choose to
reveal, and never emitting an unsigned action --- at the HTTP call, the browser
session, the keychain read, and the packet on the wire alike. The contribution is
deliberately narrow: one discipline --- sign it, seal it, cache it, bind identity
without disclosure --- that holds all the way down the stack, with each layer pinned
to an established primitive and the cache--ledger core shipped as runnable, tested
code rather than a diagram. Route discovery makes the web cheap to read~\cite{wedge};
securing the descent that executes those routes is a separate problem, and for it,
cryptography is what you needed.

\begin{thebibliography}{99}
\bibitem{wedge} L.~Tham, N.~Mac~Gregor~Garcia, J.~Hahn. \emph{Internal APIs Are All
You Need: Shadow APIs, Shared Discovery, and the Case Against Browser-First Agent
Architectures}. arXiv:2604.00694, 2026.
\bibitem{maintenance} L.~Tham. \emph{Unbrowse Maintenance Network: Proof of Indexing
and Bonded Accountability in a Shared Route Graph}. Unbrowse AI, 2026.
\bibitem{ossnotice} Unbrowse AI. \emph{Open Source Notice --- the closed/open boundary}. docs/OPEN-SOURCE-NOTICE.md, 2026.
\bibitem{x402} Coinbase. \emph{x402: An Open Protocol for Internet-Native
Payments over HTTP 402}. x402.org whitepaper, 2025. \url{https://github.com/coinbase/x402}.
\bibitem{e2e} J.~H.~Saltzer, D.~P.~Reed, D.~D.~Clark. \emph{End-to-End Arguments
in System Design}. ACM TOCS 2(4):277--288, 1984. DOI: 10.1145/357401.357402.
\bibitem{osi} ITU-T Recommendation X.200 (= ISO/IEC 7498-1), \emph{OSI Basic
Reference Model}, 1994.
\bibitem{react} S.~Yao, J.~Zhao, et al. \emph{ReAct: Synergizing Reasoning and Acting in Language Models}. ICLR 2023. arXiv:2210.03629.
\bibitem{mind2web} X.~Deng, Y.~Gu, et al. \emph{Mind2Web: Towards a Generalist Agent for the Web}. NeurIPS 2023. arXiv:2306.06070.
\bibitem{webarena} S.~Zhou, F.~F.~Xu, et al. \emph{WebArena: A Realistic Web Environment for Building Autonomous Agents}. ICLR 2024. arXiv:2307.13854.
\bibitem{seeact} B.~Zheng, B.~Gou, et al. \emph{GPT-4V(ision) is a Generalist Web Agent, if Grounded (SeeAct)}. ICML 2024. arXiv:2401.01614.
\bibitem{toolformer} T.~Schick, J.~Dwivedi-Yu, et al. \emph{Toolformer: Language Models Can Teach Themselves to Use Tools}. NeurIPS 2023. arXiv:2302.04761.
\bibitem{gorilla} S.~G.~Patil, T.~Zhang, et al. \emph{Gorilla: Large Language Model Connected with Massive APIs}. NeurIPS 2024. arXiv:2305.15334.
\bibitem{mcp} Anthropic. \emph{Model Context Protocol}. 2024. \url{https://modelcontextprotocol.io}.
\bibitem{dolevyao} D.~Dolev, A.~C.~Yao. \emph{On the Security of Public Key Protocols}. IEEE Trans. Information Theory 29(2), 1983. DOI: 10.1109/TIT.1983.1056650.
\bibitem{rfc8032} S.~Josefsson, I.~Liusvaara. \emph{Edwards-Curve Digital
Signature Algorithm (EdDSA)}. RFC 8032, IRTF CFRG, 2017.
\bibitem{erc8004} \emph{ERC-8004: Trustless Agents}. Ethereum Improvement
Proposals (Draft). \url{https://eips.ethereum.org/EIPS/eip-8004}
\bibitem{camlys} J.~Camenisch, A.~Lysyanskaya. \emph{An Efficient System for
Non-transferable Anonymous Credentials with Optional Anonymity Revocation}.
EUROCRYPT 2001, LNCS 2045. IACR ePrint 2001/019.
\bibitem{zklogin} F.~Baldimtsi, K.~Chalkias, et al. \emph{zkLogin:
Privacy-Preserving Blockchain Authentication with Existing Credentials}. ACM CCS
2024. arXiv:2401.11735.
\bibitem{semaphore} Semaphore Protocol. \emph{Semaphore: anonymous signaling on
Ethereum}. \url{https://github.com/semaphore-protocol/semaphore}
\bibitem{pedersen} T.~P.~Pedersen. \emph{Non-Interactive and
Information-Theoretic Secure Verifiable Secret Sharing}. CRYPTO 1991, LNCS 576.
DOI: 10.1007/3-540-46766-1\_9.
\bibitem{merkle} R.~C.~Merkle. \emph{A Digital Signature Based on a Conventional
Encryption Function}. CRYPTO 1987, LNCS 293. DOI: 10.1007/3-540-48184-2\_32.
\bibitem{rfc6962} B.~Laurie, A.~Langley, E.~Kasper. \emph{Certificate
Transparency}. RFC 6962, IETF, 2013.
\bibitem{bitcoin} S.~Nakamoto. \emph{Bitcoin: A Peer-to-Peer Electronic Cash
System}. 2008. \url{https://bitcoin.org/bitcoin.pdf}
\bibitem{solana} A.~Yakovenko. \emph{Solana: A New Architecture for a High
Performance Blockchain}. 2018. \url{https://solana.com/solana-whitepaper.pdf}
\bibitem{ipfs} Protocol Labs. \emph{IPFS / IPLD Merkle-DAG specifications}.
\bibitem{attention} A.~Vaswani, N.~Shazeer, et al. \emph{Attention Is All You
Need}. NeurIPS 2017. arXiv:1706.03762.
\bibitem{pagedattention} W.~Kwon, Z.~Li, et al. \emph{Efficient Memory Management
for Large Language Model Serving with PagedAttention}. SOSP 2023. arXiv:2309.06180;
\url{https://github.com/vllm-project/vllm}.
\url{https://github.com/ipfs/specs}
\bibitem{gitobjects} S.~Chacon, B.~Straub. \emph{Pro Git}, 2nd ed., \S10.2 Git
Internals --- Git Objects. \url{https://git-scm.com/book/en/v2/Git-Internals-Git-Objects}
\bibitem{rfc5116} D.~McGrew. \emph{An Interface and Algorithms for Authenticated
Encryption}. RFC 5116, 2008.
\bibitem{ocap} M.~S.~Miller. \emph{Robust Composition: Towards a Unified Approach
to Access Control and Concurrency Control}. PhD thesis, Johns Hopkins University,
2006.
\bibitem{curlimpersonate} curl-impersonate. \emph{A special build of curl that
impersonates real browsers (TLS/HTTP fingerprints)}.
\url{https://github.com/lwthiker/curl-impersonate}
\bibitem{ja3} Salesforce. \emph{JA3: TLS client fingerprinting}.
\url{https://github.com/salesforce/ja3}
\bibitem{ja4} FoxIO. \emph{JA4+ network fingerprinting suite}.
\url{https://github.com/FoxIO-LLC/ja4}
\bibitem{frost} C.~Komlo, I.~Goldberg. \emph{FROST: Flexible Round-Optimized
Schnorr Threshold Signatures}. SAC 2020, LNCS 12804. IACR ePrint 2020/852; RFC 9591.
\bibitem{ooda} J.~R.~Boyd. \emph{The Essence of Winning and Losing}, 1995; in
\emph{A Discourse on Winning and Losing}, ed. G.~Hammond, Air University Press, 2018.
\bibitem{promisetheory} M.~Burgess. \emph{An approach to understanding policy based on
autonomy and voluntary cooperation}. DSOM 2005, LNCS 3775, pp.~97--108. See also
J.~A.~Bergstra, M.~Burgess. \emph{Promise Theory: Principles and Applications}, 2014.
\bibitem{futures} B.~Liskov, L.~Shrira. \emph{Promises: Linguistic Support for Efficient
Asynchronous Procedure Calls in Distributed Systems}. PLDI 1988, pp.~260--267. Originating
the future construct: H.~C.~Baker, C.~Hewitt. \emph{The Incremental Garbage Collection of
Processes}. Proc.\ Symp.\ on AI and Programming Languages, ACM, 1977.
\end{thebibliography}

\end{document}
